\input{sol_pre.tex}
\input{q_2026-10-04_2.tex}
\begin{document}
\input{sol_head.tex}

%================= 1
\Savol{1}
Har bir qöşiluvçining 11 ga bölgandagi qoldiğini topamiz:
\[12=11+1,\quad 122=11\cdot11+1,\quad 1332=11\cdot121+1,\quad 14862=11\cdot1351+1,\quad 161061=11\cdot14641+10.\]
Qoldiqlar yiğindisi: $1+1+1+1+10=14$, \ $14=11+3$. Demak qoldiq 3.
\Javob{A) 3}

%================= 2
\Savol{2}
$\dfrac{13}{27}=\dfrac{13\cdot37}{27\cdot37}=\dfrac{481}{999}=0{,}(481)$.
\Javob{A) 0,(481)}

%================= 3
\Savol{3}
Vaqtni birinçi velosipedçi yölga çiqqan paytdan ($t=0$) hisoblaymiz. Uçinçisining tezligi $v$:
\[s_1=12t,\qquad s_2=10(t-1),\qquad s_3=v(t-2).\]
Uçinçisi ikkinçisini $t=T$ da, birinçisini $t=T+2$ da quvib yetadi:
\[v(T-2)=10(T-1),\qquad vT=12(T+2).\]
Ikkinçidan birinçisini ayiramiz: $2v=2T+34\Rightarrow v=T+17$. U holda $(T+17)T=12T+24\Rightarrow T^2+5T-24=0\Rightarrow T=3$, \ $v=20$.\\
Uçinçisi $B$ ga $70:20=3{,}5$ soatda yetadi, ya'ni $t=5{,}5$ da. Şu paytda $s_1=12\cdot5{,}5=66$, \ $s_2=10\cdot4{,}5=45$.
\[66-45=21 \text{ km}.\]
\Javob{C) 21}

%================= 4
\Savol{4}
Buyumlar jami 35 mln ga sotilgan va bu xarajatdan $40\%$ köp: xarid narxi $35:1{,}4=25$ mln.\\
Xarid narxlari $x$ va $y$:
\[\begin{cases}x+y=25\\ 1{,}25x+1{,}5y=35\end{cases}\ \Rightarrow\ 1{,}25(25-y)+1{,}5y=35\ \Rightarrow\ 0{,}25y=3{,}75\ \Rightarrow\ y=15.\]
Ikkinçi buyum xarid narxi 15 000 000 söm.
\Javob{C) 15\,000\,000}

%================= 5
\Savol{5}
$\sqrt[3]{(-5)^3}=-5$, \quad $\sqrt[3]{-5^3}=\sqrt[3]{-125}=-5$, \quad $\sqrt{(-5)^2}=|-5|=5$.
\[-5-(-5)+5=5.\]
\Javob{A) 5}

%================= 6
\Savol{6}
\[\sqrt[3]{(2\sqrt2-4)(2\sqrt2+4)}=\sqrt[3]{(2\sqrt2)^2-16}=\sqrt[3]{8-16}=\sqrt[3]{-8}=-2.\]
\Javob{D) $-2$}

%================= 7
\Savol{7}
Hammasini $\sqrt[3]{18}$ orqali ifodalaymiz ($\sqrt[3]{144}=\sqrt[3]{8\cdot18}=2\sqrt[3]{18}$):
\[30\sqrt[3]{\tfrac1{12}}=\frac{30}{\sqrt[3]{12}}=\frac{30\sqrt[3]{144}}{12}=\frac52\cdot2\sqrt[3]{18}=5\sqrt[3]{18},\qquad
\frac72\sqrt[3]{\tfrac23}=\frac72\sqrt[3]{\tfrac{18}{27}}=\frac76\sqrt[3]{18},\qquad 5\sqrt[3]{144}=10\sqrt[3]{18}.\]
\[\Bigl(5+\tfrac76+10\Bigr)\sqrt[3]{18}=\frac{97\sqrt[3]{18}}6.\]
\Javob{B) $\frac{97\sqrt[3]{18}}{6}$}

%================= 8
\Savol{8}
Çapda 11 ta, öngda 10 ta had. Toq örindagi hadlar yiğindisi $11a_{11}$ ($a_{11}$ — örta had), juft örindagilar $10a_{11}$:
\[11a_{11}=10a_{11}\ \Rightarrow\ a_{11}=0,\qquad S_{21}=21\cdot a_{11}=0.\]
\Javob{A) 0}

%================= 9
\Savol{9}
$b_4+b_5+b_6=q^3(b_1+b_2+b_3)\Rightarrow q^3=\dfrac{28}{224}=\dfrac18$, \ $q=\dfrac12$.\\
$b_1\left(1+\frac12+\frac14\right)=224\Rightarrow b_1=128$. \ $b_7=128\cdot\frac1{64}=2$, \ $b_8=1$.
\[b_7+b_8=3.\]
\Javob{C) 3}

%================= 10
\Savol{10}
Umumiy maxraj $(a+b)(b-a)=b^2-a^2$:
\[\frac{(a^2+ab+b^2)(b-a)-(a^2-ab+b^2)(a+b)}{b^2-a^2}=\frac{(b^3-a^3)-(a^3+b^3)}{b^2-a^2}=\frac{-2a^3}{b^2-a^2}=\frac{2a^3}{a^2-b^2}.\]
\Javob{C) $\frac{2a^3}{a^2-b^2}$}

%================= 11
\Savol{11}
$x=a^{\frac14}$, $y=b^{\frac14}$ deymiz. Qavs içi:
\[\frac{x-y}{x(x+y)}+\frac{x+y}{x(x-y)}-\frac{2y^2}{x(x^2-y^2)}=\frac{(x-y)^2+(x+y)^2-2y^2}{x(x^2-y^2)}=\frac{2x^2}{x(x^2-y^2)}=\frac{2x}{x^2-y^2}.\]
Uni $b^{\frac12}-a^{\frac12}=y^2-x^2$ ga köpaytiramiz: $\ \dfrac{2x}{x^2-y^2}\cdot(y^2-x^2)=-2x=-2\sqrt[4]{a}=-2\sqrt[4]2.$
\Javob{B) $-2\sqrt[4]2$}

%================= 12
\Savol{12}
$\arccos x=\frac\pi2-\arcsin x$, şuning uçun
\[\frac\pi2-\arcsin x\ge\arcsin x\ \Rightarrow\ \arcsin x\le\frac\pi4\ \Rightarrow\ x\le\frac{\sqrt2}2.\]
Aniqlaniş sohasi $[-1;1]$ ni hisobga olsak: $x\in\left[-1;\frac{\sqrt2}2\right]$.
\Javob{D)}

%================= 13
\Savol{13}
$2\sin\frac x2\cos\frac x2=\sin x$, \ $\cos^2\frac x2-\sin^2\frac x2=\cos x$:
\[\sin x\ge\cos x\ \Rightarrow\ \sqrt2\sin\!\left(x-\tfrac\pi4\right)\ge0\ \Rightarrow\ 2\pi k\le x-\tfrac\pi4\le\pi+2\pi k.\]
\[x\in\left[\frac\pi4+2\pi k;\ \frac{5\pi}4+2\pi k\right],\ k\in\mathbb Z.\]
\Javob{C)}

%================= 14
\Savol{14}
Ikkala tomonni logarifmlaymiz:
\[(2x+3)\ln7=(2x^2+3x-5)\ln5\ \Rightarrow\ 2\ln5\cdot x^2+(3\ln5-2\ln7)x-(5\ln5+3\ln7)=0.\]
Kvadrat tenglamada $a=2\ln5>0$, $c=-(5\ln5+3\ln7)<0$, ya'ni $ac<0$ $\Rightarrow$ $D>0$. Demak 2 ta ildiz.
\Javob{C) 2 ta}

%================= 15
\Savol{15}
$\log_{0{,}25}t>-1$, bu yerda $t=\log_3(4-x)>0$. Asos $0{,}25<1$ bölgani uçun işora almaşadi:
\[t<0{,}25^{-1}=4\ \Rightarrow\ 0<\log_3(4-x)<4\ \Rightarrow\ 1<4-x<81\ \Rightarrow\ -77<x<3.\]
Butun sonlar: $-76,\dots,2$ — jami $2-(-76)+1=79$ ta.
\Javob{D) 79}

%================= 16
\Savol{16}
$x\ne\pm1$. Kasrlarni qöşamiz:
\[x^2\cdot\frac{(x+1)^2+(x-1)^2}{(x^2-1)^2}=\frac{2x^2(x^2+1)}{(x^2-1)^2}=90.\]
$u=x^2$: \ $u(u+1)=45(u-1)^2\ \Rightarrow\ 44u^2-91u+45=0\ \Rightarrow\ u=\frac54$ yoki $u=\frac9{11}$.\\
Ikkalasi ham musbat va $\ne1$, har biri ikkitadan $x$ beradi: jami 4 ta ildiz.
\Javob{D) 4 ta}

%================= 17
\Savol{17}
$x\ne2$. \ $10x^2-20x=10x(x-2)$, şuning uçun
\[x^2+16=\frac{10x(x-2)}{x-2}=10x\ \Rightarrow\ x^2-10x+16=0\ \Rightarrow\ x=2\ \text{(çet)},\ x=8.\]
Yagona ildiz 8, yiğindi ham 8.
\Javob{A) 8}

%================= 18
\Savol{18}
$x=1$ da ikkala tomon 0 ga teng — $x=1$ yeçim.\\
$x\ne1$ da $|x-1|>0$ ga bölamiz:
\[\frac1{3-x}-\frac1{3-2x}\ge0\ \Rightarrow\ \frac{-x}{(3-x)(3-2x)}\ge0\ \Rightarrow\ \frac{x}{(x-3)(2x-3)}\le0.\]
Oraliqlar usuli ($0;\ 1{,}5;\ 3$ nuqtalar): $x\in(-\infty;0]\cup(1{,}5;3)$.\\
Javob: $(-\infty;0]\cup\{1\}\cup(1{,}5;3)$.
\Javob{C)}

%================= 19
\Savol{19}
Aniqlaniş sohasi: $x\ge-5$, $x\le2$, $x\ge-\frac13$ $\Rightarrow$ $x\in\left[-\frac13;2\right]$. Ikkala tomon manfiymas, kvadratga oşiramiz:
\[7+2\sqrt{(x+5)(2-x)}>3x+1\ \Rightarrow\ 2\sqrt{(x+5)(2-x)}>3x-6.\]
$x<2$ da öng tomon manfiy — tengsizlik bajariladi; $x=2$ da $0>0$ — notöğri.\\
Yeçim: $\left[-\frac13;2\right)$, butunlari $0$ va $1$ — 2 ta.
\Javob{A) 2}

%================= 20
\Savol{20}
Aniqlaniş sohasi: $x^2-6x+5\ge0$ va $x\ne1$ $\Rightarrow$ $x\in(-\infty;1)\cup[5;\infty)$. Unda $f(x)=\sqrt{x^2-6x+5}$.\\
$x^2-6x+5$ parabola $x<3$ da kamayadi, demak $f$ ham $(-\infty;1)$ da kamayadi, $[5;\infty)$ da esa ösadi.
\Javob{A) $(-\infty;1)$}

%================= 21
\Savol{21}
$\left(\dfrac{x}{9\sqrt{x^2+9}}\right)'=\dfrac{1}{(x^2+9)^{3/2}}$ (tekşiriş: hosilaning surati $\frac{x^2+9-x^2}{9}=1$). Demak
\[\int_{-4}^{4}\frac{dx}{\sqrt{(x^2+9)^3}}=\frac{x}{9\sqrt{x^2+9}}\Big|_{-4}^{4}=\frac{4}{45}-\left(-\frac4{45}\right)=\frac8{45}.\]
\Javob{A) $\frac8{45}$}

%================= 22
\Savol{22}
Grafikdagi nuqta $\left(x;\frac{x^4}8\right)$, $u=x^2\ge0$ deymiz:
\[d^2=x^2+\left(\frac{x^4}8-\frac52\right)^2=u+\left(\frac{u^2}8-\frac52\right)^2.\]
\[(d^2)'_u=1+2\left(\frac{u^2}8-\frac52\right)\frac u4=\frac{u^3-20u+16}{16}=\frac{(u-4)(u^2+4u-4)}{16}.\]
$u\ge0$ da ildizlar: $u=4$ va $u=2\sqrt2-2$. Qiymatlar: $u=4$: $d^2=4+\frac14=\frac{17}4$; \ $u=2\sqrt2-2\approx0{,}83$: $d^2\approx6{,}7$; \ $u=0$: $d^2=6{,}25$.
\[d_{\min}=\sqrt{\tfrac{17}4}=\frac{\sqrt{17}}2.\]
\Javob{A) $\frac{\sqrt{17}}2$}

%================= 23
\Savol{23}
\fig{\begin{tikzpicture}[scale=0.32]
\draw[thick] (0,0) circle (8);
\coordinate (O) at (0,0); \coordinate (A) at (-5.657,-5.657); \coordinate (B) at (5.657,-5.657);
\draw[thick] (A)--(B);
\draw[axs] (O)--(A); \draw[axs] (O)--(B);
\draw[aux] ($(O)!0.9cm!(A)$)--($(O)!0.9cm!(A)+(O)!0.9cm!(B)$)--($(O)!0.9cm!(B)$);
\fill (O) circle (4pt);
\node[above] at (O){$O$}; \node[below left] at (A){$A$}; \node[below right] at (B){$B$};
\node[aux,above left] at ($(O)!0.5!(A)$){$8$}; \node[aux,above right] at ($(O)!0.5!(B)$){$8$};
\end{tikzpicture}}
$\angle AOB=90^\circ$, $OA=OB=8$ (\textcolor{aux}{qizil} radiuslar). $\triangle AOB$ — teng yonli töğri burçakli:
\[AB=\sqrt{8^2+8^2}=8\sqrt2.\]
\Javob{A) $8\sqrt2$}

%================= 24
\Savol{24}
$Ox$ böylab 3 birlik öngga va $Oy$ böylab 2 birlik yuqoriga: \ $g(x)=f(x-3)+2$.
\[g(2)=f(-1)+2=\bigl(-1+1+9+8\bigr)+2=19.\]
\Javob{B) 19}

%================= 25
\Savol{25}
\fig{\begin{tikzpicture}[scale=1.0]
\coordinate (A) at (0,0); \coordinate (C) at (5,0); \coordinate (B) at (3.943,3.943);
\coordinate (D) at (1.057,1.057); \coordinate (E) at (2.887,0);
\coordinate (M) at ($(D)!0.5!(C)$);
\draw[thick] (A)--(B)--(C)--cycle; \draw[thick] (B)--(E); \draw[thick] (D)--(C); \draw[thick] (D)--(E);
\draw[axs] (E)--(C);
\draw[aux] ($(M)!0.22cm!(C)$)--($($(M)!0.22cm!(C)$)!1!90:(M)$)--($(M)!0.22cm!(B)$);
\fill (M) circle (1.3pt);
\pic[draw,angle radius=0.6cm,"\scriptsize$45^\circ$",angle eccentricity=1.45] {angle=C--A--B};
\pic[draw,angle radius=0.9cm,"\scriptsize$30^\circ$",angle eccentricity=1.35] {angle=A--B--E};
\pic[draw,angle radius=1.3cm,"\scriptsize$30^\circ$",angle eccentricity=1.25] {angle=E--B--C};
\pic[draw,angle radius=0.7cm,"\scriptsize$60^\circ$",angle eccentricity=1.4] {angle=B--C--D};
\pic[draw,angle radius=1.3cm,"\scriptsize$15^\circ$",angle eccentricity=1.25] {angle=D--C--A};
\node[left] at (A){$A$}; \node[above] at (B){$B$}; \node[right] at (C){$C$}; \node[above left] at (D){$D$}; \node[below] at (E){$E$};
\node[aux,below] at ($(M)+(0,-0.05)$){\scriptsize$M$};
\end{tikzpicture}}
$\angle ACB=15^\circ+60^\circ=75^\circ$, \ $\angle ABC=180^\circ-45^\circ-75^\circ=60^\circ$, \ $\angle ABE=\angle EBC=30^\circ$.\\
$\triangle DBC$ da $\angle DBC=\angle BCD=60^\circ$ — u teng tomonli. Bissektrisa $BE$ teng tomonli uçburçakda $DC$ ning örta perpendikulyari böladi ($M$ — $DC$ örtasi), şuning uçun $ED=EC$ (\textcolor{aux}{qizil}) va $\angle EDC=\angle ECD=15^\circ$.\\
$\triangle DME$ ($\angle M=90^\circ$): \ $\angle DEB=\angle DEM=90^\circ-15^\circ=75^\circ$.
\Javob{D) $75^\circ$}

%================= 26
\Savol{26}
\fig{\begin{tikzpicture}[scale=1.05]
\coordinate (A) at (0,0); \coordinate (M) at (2.8,0); \coordinate (H) at (4.2,0); \coordinate (C) at (5.6,0); \coordinate (B) at (4.2,3.85);
\draw[thick] (A)--(B)--(C)--cycle; \draw[thick] (B)--(M); \draw[ax] (B)--(H);
\draw[aux] (4.2,0.22)--(4.42,0.22)--(4.42,0);
\draw[axs] (B)--(M); \draw[axs] (B)--(C);
\pic[draw,angle radius=0.8cm,"\scriptsize$40^\circ$",angle eccentricity=1.4] {angle=M--B--C};
\pic[draw,angle radius=0.45cm] {angle=B--M--A};
\node[below] at (A){$A$}; \node[below] at (M){$M$}; \node[below] at (H){$H$}; \node[below] at (C){$C$}; \node[above] at (B){$B$};
\node[below] at (1.4,-0.3){\small$28$}; \node[below] at (3.5,-0.3){\small$14$}; \node[below] at (4.9,-0.3){\small$14$};
\end{tikzpicture}}
$M$ — $AC$ ning örtasi: $AM=MC=28$. $MH=14$ bölgani uçun $HC=28-14=14=MH$ ($H$ nuqta $M$ va $C$ orasida; aks holda $\angle AMB$ ötkir bölib, javoblarda yöq).\\
$BH$ — $MC$ ga perpendikulyar va uni teng ikkiga böladi, demak $BM=BC$ (\textcolor{aux}{qizil}) — $\triangle MBC$ teng yonli:
\[\angle BMC=\angle BCM=\frac{180^\circ-40^\circ}2=70^\circ,\qquad \angle AMB=180^\circ-70^\circ=110^\circ.\]
\Javob{D) $110^\circ$}

%================= 27
\newpage
\Savol{27}
\fig{\begin{tikzpicture}[scale=0.95]
\coordinate (A) at (0,0); \coordinate (B) at (0.6,3.6); \coordinate (C) at (5.6,3.6); \coordinate (D) at (5,0);
\coordinate (E) at ($(A)!0.25!(B)$); \coordinate (F) at ($(A)!0.75!(B)$);
\coordinate (M) at ($(C)!1/3!(D)$); \coordinate (N) at ($(C)!2/3!(D)$);
\coordinate (O) at (intersection of F--N and E--M);
\fill[msblue!25] (E)--(F)--(O)--cycle; \fill[msblue!25] (M)--(N)--(O)--cycle;
\draw[thick] (A)--(B)--(C)--(D)--cycle; \draw (F)--(N); \draw (E)--(M);
\coordinate (P) at ($(A)!(O)!(B)$); \coordinate (Q) at ($(D)!(O)!(C)$);
\draw[ax] (O)--(P); \draw[ax] (O)--(Q);
\node[aux,above] at ($(O)!0.5!(P)+(0,0.05)$){\scriptsize$3k$}; \node[aux,above] at ($(O)!0.5!(Q)+(0,0.05)$){\scriptsize$2k$};
\foreach \p in {A,B,C,D,E,F,M,N,O} \fill (\p) circle (1.2pt);
\node[below left] at (A){$A$}; \node[above left] at (B){$B$}; \node[above right] at (C){$C$}; \node[below right] at (D){$D$};
\node[left] at (E){$E$}; \node[left] at (F){$F$}; \node[right] at (M){$M$}; \node[right] at (N){$N$}; \node[below] at ($(O)+(0,-0.08)$){$O$};
\node[left,font=\small] at ($(A)!0.5!(E)$){$x$}; \node[left,font=\small] at ($(E)!0.5!(F)+(0,0.25)$){$2x$}; \node[left,font=\small] at ($(F)!0.5!(B)$){$x$};
\end{tikzpicture}}
$AE=BF=x$, $EF=2x$, $AB=4x$; \ $CM=MN=ND=\frac{4x}3$.\\
$EF\parallel MN$ bölgani uçun $\triangle EFO\sim\triangle MNO$, öxşaşlik koeffitsiyenti $\dfrac{EF}{MN}=\dfrac{2x}{4x/3}=\dfrac32$.\\
Demak $O$ dan $AB$ va $CD$ gaça masofalar ham $3:2$ (\textcolor{aux}{qizil}). Parallelogramm balandligi ($AB$ va $CD$ orasida) $h$ bölsa, ular $\frac35h$ va $\frac25h$.
\[S_{EFO}+S_{MNO}=\frac12\cdot2x\cdot\frac35h+\frac12\cdot\frac{4x}3\cdot\frac25h=\frac{3xh}5+\frac{4xh}{15}=\frac{13xh}{15},\qquad S_{ABCD}=4xh.\]
\[\frac{13xh/15}{4xh}=\frac{13}{60}.\]
\Javob{B) $\frac{13}{60}$}

%================= 28
\Savol{28}
\fig{\begin{tikzpicture}[scale=0.5]
\def\s{1.732}
\coordinate (L) at (-4.5,0); \coordinate (R) at (4.5,0); \coordinate (T) at (0,7.794);
\foreach \cx/\cy in {-1.5/1.732, 1.5/1.732, 0/4.330}{
 \draw[thick,fill=msblue!25] ({\cx+\s*cos(30)},{\cy+\s*sin(30)}) \foreach \k in {1,...,5}{ -- ({\cx+\s*cos(30+60*\k)},{\cy+\s*sin(30+60*\k)}) } -- cycle;}
\draw[thick] (L)--(R)--(T)--cycle;
\coordinate (O1) at (-1.5,1.732); \coordinate (V) at (-3,2.598); \coordinate (W) at (-1.5,0);
\draw[ax] (L)--(O1); \draw[axs] (O1)--(V); \draw[axs] (O1)--(W);
\fill[aux] (O1) circle (3pt); \fill (V) circle (3pt); \fill (W) circle (3pt);
\node[aux,font=\small] at (-1.2,0.75){$a$}; \node[aux,font=\small] at (-2.1,2.6){$a$};
\node[left] at (V){$V$}; \node[below] at (W){$W$}; \node[aux,right] at (O1){$O_1$};
\node[below left] at (L){$A$};
\node[aux,font=\scriptsize] at (-3.25,0.2){$30^\circ$};
\draw[|<->|] (-4.5,-1.3)--node[lbl]{$a\sqrt3$}(-1.5,-1.3); \draw[|<->|] (-1.5,-1.3)--node[lbl]{$a\sqrt3$}(1.5,-1.3); \draw[|<->|] (1.5,-1.3)--node[lbl]{$a\sqrt3$}(4.5,-1.3);
\end{tikzpicture}}
Oltiburçak tomoni $a$; uning markazidan har bir uçigaça masofa $a$, qarama-qarşi parallel tomonlar orasidagi masofa $a\sqrt3$.\\
Çap pastki oltiburçak (markazi $O_1$) asosga $W$ uçi bilan, çap tomonga $V$ uçi bilan tegadi. $O_1W$ asosga, $O_1V$ çap tomonga perpendikulyar (\textcolor{aux}{qizil}), ikkalasi ham $a$ ga teng. Demak $O_1$ — $A$ burçak bissektrisasida: $\angle O_1AW=30^\circ$,
\[AW=\frac{O_1W}{\tg30^\circ}=a\sqrt3.\]
Ikki pastki oltiburçak umumiy vertikal tomonga ega, markazlari orasidagi masofa $a\sqrt3$. Simmetriya böyiça öngda ham $a\sqrt3$:
\[9=a\sqrt3+a\sqrt3+a\sqrt3=3\sqrt3\,a\ \Rightarrow\ a=\sqrt3.\]
\[S=3\cdot\frac{3\sqrt3}2a^2=3\cdot\frac{3\sqrt3}{2}\cdot3=\frac{27\sqrt3}2.\]
\Javob{D) $\frac{27\sqrt3}2$}

%================= 29
\Savol{29}
\fig{\begin{tikzpicture}[scale=0.9]
\draw[->] (-0.5,0)--(4.8,0) node[right]{$x$}; \draw[->] (0,-0.5)--(0,3.8) node[above]{$y$};
\coordinate (A) at (0,0); \coordinate (B) at (1,2); \coordinate (C) at (3,3); \coordinate (D) at (4,0); \coordinate (E) at (3.375,1.875);
\fill[aux!15] (A)--(E)--(D)--cycle;
\draw[thick] (A)--(B)--(C)--(D)--cycle; \draw[axs] (A)--(E);
\draw[ax] (E)--(3.375,0);
\foreach \p in {A,B,C,D,E} \fill (\p) circle (1.5pt);
\node[below left] at (A){$A$}; \node[left] at (B){$B(1;2)$}; \node[above] at (C){$C(3;3)$}; \node[below] at (D){$D(4;0)$}; \node[right] at (E){$E$};
\node[aux,font=\small] at (3.6,0.9){$y_E$};
\end{tikzpicture}}
Törtburçak yuzini bölaklab hisoblaymiz: $B$ va $C$ dan $Ox$ ga perpendikulyarlar tuşiramiz ($x=1$ va $x=3$ da):
\[S=\underbrace{\tfrac12\cdot1\cdot2}_{\text{uçburçak}}+\underbrace{\tfrac{2+3}2\cdot2}_{\text{trapetsiya}}+\underbrace{\tfrac12\cdot1\cdot3}_{\text{uçburçak}}=1+5+1{,}5=7{,}5.\]
$AE$ yuzni teng ikkiga böladi: $S_{AED}=3{,}75$ (\textcolor{aux}{qizil} uçburçak). Uning asosi $AD=4$, balandligi $y_E$:
\[\frac12\cdot4\cdot y_E=3{,}75\ \Rightarrow\ y_E=\frac{15}8.\]
$CD$ töğri çiziği: $y=-3(x-4)$, \ $\frac{15}8=-3(x_E-4)\Rightarrow x_E=4-\frac58=\frac{27}8$.
\Javob{C) $\left(\frac{27}8;\frac{15}8\right)$}

%================= 30
\Savol{30}
\fig{\begin{tikzpicture}[scale=0.85]
\coordinate (A) at (0,0); \coordinate (B) at (4,0); \coordinate (C) at (5.4,1.4); \coordinate (D) at (1.4,1.4);
\coordinate (O) at (2.7,0.7); \coordinate (S) at (2.7,3.6); \coordinate (K) at ($(A)!0.5!(B)$);
\draw[thick] (A)--(B)--(C); \draw[dashed] (C)--(D)--(A); \draw[thick] (S)--(A); \draw[thick] (S)--(B); \draw[thick] (S)--(C); \draw[dashed] (S)--(D);
\draw[ax] (S)--(O); \draw[ax] (O)--(K); \draw[axs] (S)--(K);
\pic[draw=aux,angle radius=0.45cm,"\scriptsize$45^\circ$",angle eccentricity=1.7] {angle=O--K--S};
\foreach \p in {S,O,K} \fill (\p) circle (1.3pt);
\node[above] at (S){$S$}; \node[right] at (O){$O$}; \node[below] at (K){$K$};
\node[below left] at (A){$A$}; \node[below right] at (B){$B$}; \node[right] at (C){$C$}; \node[above left] at (D){$D$};
\end{tikzpicture}}
Asos tomoni $a$, apofema $SK=l$. $\triangle SOK$ da $\angle SKO=45^\circ$, $OK=\frac a2$ (\textcolor{aux}{qizil}):
\[l=\frac{OK}{\cos45^\circ}=\frac{a}{\sqrt2},\qquad H=SO=OK=\frac a2.\]
Yon sirt: $4\cdot\frac12al=\frac{2a^2}{\sqrt2}=\sqrt2a^2=144\sqrt2\Rightarrow a=12$, \ $H=6$.
\[V=\frac13a^2H=\frac13\cdot144\cdot6=288.\]
\Javob{B) 288}

%================= 31
\Savol{31}
\fig{\begin{tikzpicture}[x=0.32mm,y=-0.32mm]
\def\ra{42}
\begin{scope}\clip (483,763) circle (\ra); \fill[aux!20] (422,763) circle (\ra); \fill[aux!20] (452,808) circle (\ra);\end{scope}
\draw (346,696) rectangle (559,875);
\draw (422,763) circle (\ra) (483,763) circle (\ra) (452,808) circle (\ra);
\node at (358,684){$U$}; \node at (405,711){$A$}; \node at (508,711){$B$}; \node at (492,858){$C$};
\foreach \x/\y/\t in {370/747/t,370/785/r,370/820/o,391/849/s,405/758/a,424/740/b,407/773/c,487/744/d,504/756/e,504/773/f,428/780/u,440/792/x,430/830/k,451/837/n,469/821/m}{\node at (\x,\y){\small$\t$};}
\foreach \x/\y/\t in {453/745/p,453/773/q,473/785/l,479/797/v}{\node[aux] at (\x,\y){\small$\boldsymbol{\t}$};}
\end{tikzpicture}}
$B'$ — $B$ dan taşqaridagi elementlar. Biror töplamdan $B'$ ni ayiriş — undagi $B$ ga kirmaydigan elementlarni olib taşlaş:
\[(A\cup C)\setminus B'=(A\cup C)\cap B.\]
Bu — $B$ ning $A$ yoki $C$ bilan umumiy qismi (\textcolor{aux}{qizil} soha): $\{p;\ q;\ l;\ v\}$ — 4 ta element.
\Javob{A) 4}

%================= 32
\Savol{32}
$\overline{abca}$: $a$ — birinçi raqam, $a\ne0$: 9 ta tanlov; $b\ne a$: 9 ta; $c\ne a,\ c\ne b$: 8 ta. Oxirgi raqam $a$ ga teng — tanlov yöq.
\[9\cdot9\cdot8=648.\]
\Javob{C) 648}

%================= 33-35
\newpage
\Savol{33}
Silindr öq kesimi kvadrat: $h=2r$. Kub qirrasi $a$. \ $\pi\approx3$.\\
Silindr: $V_s=\pi r^2\cdot2r=6r^3$, \ yon sirt $2\pi r\cdot2r=12r^2$, \ töla sirt $12r^2+2\pi r^2=18r^2$, \ asos yuzi $3r^2$.\\
Kub: $V_k=a^3$, \ yon sirt $4a^2$, \ töla sirt $6a^2$, \ yoq yuzi $a^2$.

\textbf{33.} $6r^3=a^3$. \ Nisbat $\dfrac{18r^2}{6a^2}=\dfrac{3r^2}{a^2}$, uning kubi:
\[\frac{27r^6}{a^6}=\frac{27r^6}{(6r^3)^2}=\frac{27}{36}=\frac34.\]
\textbf{34.} $12r^2=4a^2\Rightarrow a^2=3r^2$. \ $\dfrac{a^2}{3r^2}=1$.

\textbf{35.} $a=r$: \ $\dfrac{a^3}{6r^3}=\dfrac16$.
\Javob{33 – A; \ 34 – E; \ 35 – C}

%================= 36
\newpage
\Savol{36}
$x=0$ ildiz faqat $k^2=0$, ya'ni $k=0$ da. Avval $k\ne0$ deb olamiz va tenglamani $x^2$ ga bölamiz:
\[\Bigl(x+\frac kx+3\Bigr)\Bigl(x+\frac kx+9\Bigr)+8=0.\]
$t=x+\dfrac kx$: \quad $(t+3)(t+9)+8=t^2+12t+35=0\ \Rightarrow\ t=-5$ yoki $t=-7$.
\[x+\frac kx=-5\iff x^2+5x+k=0\quad(D_1=25-4k),\qquad x+\frac kx=-7\iff x^2+7x+k=0\quad(D_2=49-4k).\]
Bu ikki tenglamaning umumiy ildizi bölişi uçun $5x=7x$, ya'ni $x=0$ kerak — $k\ne0$ da bunday emas. Demak ildizlar soni:
\begin{itemize}\setlength{\itemsep}{1pt}
\item $k<\frac{25}4$ ($k\ne0$): $D_1>0,\ D_2>0$ — \textbf{4 ta} ildiz;
\item $k=\frac{25}4$: $1+2=3$ ta; \quad $\frac{25}4<k<\frac{49}4$: $D_1<0,\ D_2>0$ — \textbf{2 ta}; \quad $k=\frac{49}4$: 1 ta; \quad $k>\frac{49}4$: yöq.
\item $k=0$: $x^2(x^2+12x+35)=0\Rightarrow x\in\{0;-5;-7\}$ — 3 ta.
\end{itemize}
a) 2 ta ildiz: $6{,}25<k<12{,}25$, butunlari $7,8,\dots,12$: \ $7+8+9+10+11+12=57$.

b) 4 ta ildiz: $k<6{,}25$, $k\ne0$ — eng katta butun qiymat $k=6$.
\Javob{a) 57; \quad b) 6}

%================= 37
\newpage
\Savol{37}
Köpaytmani yiğindiga keltiramiz:
\[8\cos\frac{3x}4\cos\frac x4=4\Bigl(\cos x+\cos\frac x2\Bigr).\]
Tenglama: $4\cos x+4\cos\frac x2-6\cos x+1=0\ \Rightarrow\ -2\cos x+4\cos\frac x2+1=0$.\\
$c=\cos\frac x2$, \ $\cos x=2c^2-1$:
\[-2(2c^2-1)+4c+1=0\ \Rightarrow\ 4c^2-4c-3=0\ \Rightarrow\ c=\frac32\ \text{(mos emas)}\ \ \text{yoki}\ \ c=-\frac12.\]
\[\cos\frac x2=-\frac12\ \Rightarrow\ \frac x2=\pm\frac{2\pi}3+2\pi n\ \Rightarrow\ x=\pm\frac{4\pi}3+4\pi n,\ n\in\mathbb Z.\]
a) Eng kiçik musbat ildiz: $x=\dfrac{4\pi}3$.

b) $(-2\pi;\,2\pi)$ oraliqda: $\dfrac{4\pi}3$ va $-\dfrac{4\pi}3$. Keyingilari $\pm\frac{4\pi}3\mp4\pi=\mp\frac{8\pi}3$ — oraliqdan taşqarida. Jami 2 ta.
\Javob{a) $\frac{4\pi}3$; \quad b) 2}

%================= 38
\newpage
\Savol{38}
Tenglikda $x$ ni $-x$, $\frac1x$, $-\frac1x$ bilan almaştiramiz. Belgilaymiz: $A=f(x)$, $B=f(-x)$, $C=f\!\left(\frac1x\right)$, $D=f\!\left(-\frac1x\right)$:
\[\begin{aligned}
&(1)\ \ 3B+A+C=x, &&(2)\ \ 3A+B+D=-x,\\
&(3)\ \ 3D+C+A=\tfrac1x, &&(4)\ \ 3C+D+B=-\tfrac1x.
\end{aligned}\]
$(1)+(2)$: $4(A+B)+(C+D)=0$; \ $(3)+(4)$: $(A+B)+4(C+D)=0$ \ $\Rightarrow$ \ $A+B=0$, $C+D=0$, ya'ni $f$ — toq funksiya.\\
Unda (1): $-3A+A+C=x\Rightarrow C=x+2A$; \ (3): $-3C+C+A=\frac1x\Rightarrow A-2C=\frac1x$:
\[A-2x-4A=\frac1x\ \Rightarrow\ f(x)=A=-\frac{2x^2+1}{3x}.\]
a) $f(1)=-\dfrac{3}{3}=-1$.

b) $f(x)=1$: \ $-(2x^2+1)=3x\Rightarrow2x^2+3x+1=0\Rightarrow x_1=-1,\ x_2=-\frac12$.
\[\frac1{x_1^3}+\frac1{x_2^3}=-1+(-8)=-9.\]
\Javob{a) $-1$; \quad b) $-9$}

%================= 39
\newpage
\Savol{39}
Tengliklarni $x$ böyiça differensiallaymiz (murakkab funksiya hosilasi):
\[3f'(3x+1)=3x^2-6x+12\ \Rightarrow\ f'(3x+1)=x^2-2x+4,\]
\[-2g'(1-2x)=2-2x\ \Rightarrow\ g'(1-2x)=x-1.\]
a) $3x+1=4\Rightarrow x=1$: $f'(4)=1-2+4=3$. \quad $1-2x=-3\Rightarrow x=2$: $g'(-3)=1$.
\[\frac{f'(4)}{g'(-3)}=3.\]
b) $1-2x=-7\Rightarrow x=4$: \ $g(-7)=2\cdot4-16=-8$.\\
$3x+1=-8\Rightarrow x=-3$: \ $f'(-8)=9+6+4=19$.
\Javob{a) 3; \quad b) 19}

%================= 40
\newpage
\Savol{40}
\fig{\begin{tikzpicture}[xscale=0.45,yscale=0.14]
\fill[okgreen!30,domain=-4:3,samples=60] plot (\x,{-\x*\x+2*\x+8}) -- plot[domain=3:-4] (\x,{3*\x-4});
\draw[->] (-5.2,0)--(5.4,0) node[right]{$x$}; \draw[->] (0,-18)--(0,11) node[above]{$y$};
\draw[thick,domain=-4.5:4.5,samples=60] plot (\x,{-\x*\x+2*\x+8});
\draw[thick,domain=-4.6:4.3] plot (\x,{3*\x-4});
\draw[ax] (3,0)--(3,5); \draw[ax] (0,5)--(3,5);
\draw[ax] (-4,0)--(-4,-16); \draw[ax] (0,-16)--(-4,-16);
\foreach \p in {(-2,0),(4,0),(0,8),(3,5),(-4,-16)} \fill \p circle (2.2pt);
\node[below right] at (-2,0){$-2$}; \node[below right] at (4,0){$4$}; \node[left] at (0,8.6){$8$};
\node[right] at (3.1,6){$A(3;5)$}; \node[below] at (3,0){\color{aux}$3$}; \node[right] at (0,5){\color{aux}$5$};
\node[left] at (-4.1,-16){$(-4;-16)$}; \node[right] at (0,-16){\color{aux}$-16$};
\end{tikzpicture}}
\textbf{Parabola} $(-2;0)$, $(4;0)$ dan ötadi: $y=a(x+2)(x-4)$; \ $(0;8)$: $-8a=8\Rightarrow a=-1$:
\[y=-(x+2)(x-4)=-x^2+2x+8,\qquad a=-1,\ b=2,\ c=8.\]
\textbf{Kesişiş}: $-x^2+2x+8=3x-4\Rightarrow x^2+x-12=0\Rightarrow x=3$ yoki $x=-4$. \ $m=3>0$, $n=3\cdot3-4=5$.

a) \[\frac{a+b+c}{m+n}=\frac{-1+2+8}{3+5}=\frac98.\]
b) Böyalgan soha $[-4;3]$ da parabola ostida, töğri çiziq ustida:
\[S=\int_{-4}^{3}\bigl(-x^2+2x+8-(3x-4)\bigr)dx=\int_{-4}^{3}(-x^2-x+12)\,dx=\left(-\frac{x^3}3-\frac{x^2}2+12x\right)\Big|_{-4}^{3}\]
\[=\left(-9-\frac92+36\right)-\left(\frac{64}3-8-48\right)=\frac{45}2+\frac{104}3=\frac{343}6.\]
\Javob{a) $\frac98$; \quad b) $\frac{343}6$}

%================= 41
\newpage
\Savol{41}
\fig{\begin{tikzpicture}[scale=0.27]
\coordinate (A) at (0,0); \coordinate (B) at (8,0); \coordinate (C) at (6.899,11.949);
\coordinate (D) at (4,0); \coordinate (E) at ($(B)!0.5!(C)$); \coordinate (O) at (4,5.657);
\coordinate (F) at (9.610,9.717); \coordinate (G) at (3.288,-1.234);
\draw[thick] (O) circle (6.928);
\draw[thick] (A)--(B)--(C)--cycle; \draw[thick] (G)--(F);
\draw[axs] ($(A)+(0,0)$)--(D); \draw[axs] (B)--(E);
\foreach \p in {A,B,C,D,E,F,G} \fill (\p) circle (4pt);
\node[left] at (A){$A$}; \node[right] at (B){$B$}; \node[above] at (C){$C$}; \node at (4.6,-0.6){$D$};
\node[left] at (E){$E$}; \node[right] at (F){$F$}; \node[below] at (G){$G$};
\node[aux] at (2.8,0.85){\small$4$}; \node[aux] at (6.2,0.85){\small$4$};
\node[aux,right] at ($(B)!0.5!(E)+(0.1,0)$){\small$6$}; \node[right] at ($(E)!0.5!(C)+(0.15,0.2)$){\small$6$};
\pic[draw,angle radius=0.5cm,"\scriptsize$60^\circ$",angle eccentricity=1.7] {angle=B--A--C};
\end{tikzpicture}}
Sinuslar teoremasi ($2R=8\sqrt3$):
\[BC=2R\sin60^\circ=8\sqrt3\cdot\frac{\sqrt3}2=12,\qquad \sin C=\frac{AB}{2R}=\frac8{8\sqrt3}=\frac1{\sqrt3},\ \ \cos C=\sqrt{\tfrac23}.\]
$\sin B=\sin(120^\circ-C)=\frac{\sqrt3}2\cdot\sqrt{\frac23}+\frac12\cdot\frac1{\sqrt3}=\frac{\sqrt2}2+\frac1{2\sqrt3}$, \quad $AC=2R\sin B=4\sqrt6+4$.

a) $DE$ — örta çiziq: $DE=\dfrac{AC}2=2+2\sqrt6$.

b) Kesişuvçi vatarlar xossasi (\textcolor{aux}{qizil}: $AD=DB=4$, $BE=EC=6$). $GD=u$, $EF=w$, $DE=d$:
\[GD\cdot DF=AD\cdot DB:\ u(d+w)=16;\qquad FE\cdot EG=BE\cdot EC:\ w(d+u)=36.\]
Ayiramiz: $d(w-u)=20\Rightarrow w=u+\frac{20}d$. \ $\frac{20}{d}=\frac{20}{2(1+\sqrt6)}=2\sqrt6-2$, \ $k=d+\frac{20}d=4\sqrt6$.\\
Birinçi tenglama: $u(u+k)=16\Rightarrow\left(u+\frac k2\right)^2=16+\frac{k^2}4$. \ $GF=u+d+w=2u+k=2\left(u+\frac k2\right)$:
\[GF=\sqrt{64+k^2}=\sqrt{64+96}=\sqrt{160}=4\sqrt{10}.\]
\Javob{a) $2+2\sqrt6$; \quad b) $4\sqrt{10}$}

%================= 42
\newpage
\Savol{42}
\fig{\begin{tikzpicture}[scale=0.62]
\def\s{4.123}
\coordinate (A) at (0,0); \coordinate (B) at (0,\s); \coordinate (C) at (\s,\s); \coordinate (D) at (\s,0); \coordinate (E) at (1.213,0.728);
\coordinate (F) at (4.851,2.910);
\fill[msblue!20] (A)--(C)--(E)--cycle;
\draw[thick] (A)--(B)--(C)--(D)--cycle; \draw[thick] (A)--(C); \draw[thick] (A)--(E)--(C); \draw[thick] (E)--(D);
\draw[axs] (D)--(F)--(C); \draw[ax] (E)--(F);
\pic[draw=aux,angle radius=0.25cm]{right angle=F--D--E}; \pic[draw=aux,angle radius=0.2cm]{right angle=C--F--E};
\fill (E) circle (1.5pt); \fill[aux] (F) circle (1.5pt);
\node[below left] at (A){$A$}; \node[above left] at (B){$B$}; \node[above right] at (C){$C$}; \node[below right] at (D){$D$};
\node[below right] at (E){$E$}; \node[aux,right] at (F){$E'$};
\node[aux,right] at ($(D)!0.5!(F)+(0.05,-0.1)$){\scriptsize$3$}; \node[aux,right] at ($(C)!0.5!(F)+(0.08,0.05)$){\scriptsize$\sqrt2$};
\end{tikzpicture}}
Kvadrat tomoni $s$. $\triangle ADE$ ni $D$ atrofida $90^\circ$ ga buramiz: $A\to C$, $E\to E'$ (\textcolor{aux}{qizil}).
\[DE'=DE=3,\qquad CE'=AE=\sqrt2,\qquad \angle EDE'=90^\circ.\]
$\triangle EDE'$ — teng yonli töğri burçakli: $EE'=3\sqrt2$, \ $\angle DE'E=45^\circ$.\\
$\triangle EE'C$ da $EE'^2+CE'^2=18+2=20=EC^2$ — Pifagor teoremasiga teskari teoremaga köra $\angle EE'C=90^\circ$. Demak
\[\angle DE'C=\angle DE'E+\angle EE'C=45^\circ+90^\circ=135^\circ.\]
a) $\triangle DE'C$ da kosinuslar teoremasi:
\[s^2=DC^2=3^2+(\sqrt2)^2-2\cdot3\cdot\sqrt2\cdot\cos135^\circ=9+2+6=17\ \Rightarrow\ S_{ABCD}=17.\]
b) $S_{ACE}=S_{ACD}-S_{AED}-S_{DEC}$. Uçburçak yuzini tomonlari orqali (Geron formulasining $16S^2=2(a^2b^2+b^2c^2+c^2a^2)-(a^4+b^4+c^4)$ körinişi):
\[\triangle AED\ (2;\,9;\,17):\ \ 16S^2=2(18+153+34)-(4+81+289)=36\ \Rightarrow\ S_{AED}=\frac32,\]
\[\triangle DEC\ (9;\,20;\,17):\ 16S^2=2(180+340+153)-(81+400+289)=576\ \Rightarrow\ S_{DEC}=6.\]
(Qavslarda tomonlar kvadratlari.) \quad $S_{ACE}=\dfrac{17}2-\dfrac32-6=1.$
\Javob{a) 17; \quad b) 1}

%================= 43
\newpage
\Savol{43}
\fig{\begin{tikzpicture}[scale=0.19]
\coordinate (A) at (0,0); \coordinate (C) at (27,0); \coordinate (O1) at (10,0); \coordinate (B) at (7.714,10.286);
\coordinate (O2) at (13.5,-2.09);
\draw[thick] (O2) circle (13.661); \draw[thick] (O1) circle (8);
\draw[thick] (A)--(B)--(C)--cycle;
\coordinate (P) at ($(A)!(O1)!(B)$); \coordinate (Q) at ($(B)!(O1)!(C)$);
\draw[axs] (O1)--(P); \draw[axs] (O1)--(Q); \draw[ax] (B)--(O1);
\foreach \p in {A,B,C,O1,P,Q} \fill (\p) circle (5pt);
\fill (2,0) circle (5pt); \fill (18,0) circle (5pt);
\node[left] at (A){$A$}; \node[above] at (B){$B$}; \node[right] at (C){$C$};
\node[below] at ($(O1)+(0.4,-0.1)$){$O_1$}; \node at (3.0,-1.0){\small$K$}; \node at (16.9,-1.0){\small$L$};
\node[aux,left] at ($(O1)!0.5!(P)$){\small$8$}; \node[aux,above] at ($(O1)!0.55!(Q)$){\small$8$};
\node[above] at (22.5,0){\small$9$};
\end{tikzpicture}}
Kiçik aylana radiusi $r=\frac{KL}2=8$. $AO_1=2+8=10$, \ $O_1C=8+9=17$, \ $AC=27$.\\
Aylana $AB$ va $BC$ ga urinadi, demak $O_1$ undan teng uzoqlikda: $BO_1$ — $B$ burçak bissektrisasi (\textcolor{aux}{qizil}).

a) Bissektrisa xossasi: \ $\dfrac{AB}{BC}=\dfrac{AO_1}{O_1C}=\dfrac{10}{17}$.

b) $O_1$ dan $AB$ va $BC$ gaça masofa 8:
\[\sin A=\frac8{10}=\frac45,\ \cos A=\frac35;\qquad \sin C=\frac8{17},\ \cos C=\frac{15}{17}.\]
\[\sin B=\sin(A+C)=\frac45\cdot\frac{15}{17}+\frac35\cdot\frac8{17}=\frac{84}{85},\qquad R=\frac{AC}{2\sin B}=\frac{27\cdot85}{2\cdot84}=\frac{765}{56}.\]
\Javob{a) $\frac{10}{17}$; \quad b) $\frac{765}{56}$}

%================= 44
\newpage
\Savol{44}
\fig{\begin{tikzpicture}[scale=0.3]
\coordinate (O) at (0,0); \coordinate (O1) at (0,12); \coordinate (A) at (10.392,0); \coordinate (A1) at (3.464,12); \coordinate (Q) at (0,2); \coordinate (P) at (3.464,0);
\draw[thick] (O)--(A)--(A1)--(O1); \draw[ax] (O)--(O1);
\draw[ax] (A1)--(P);
\draw[axs] (Q)--(A); \draw[axs] (Q)--(A1);
\pic[draw,angle radius=0.5cm,"\scriptsize$60^\circ$",angle eccentricity=1.7] {angle=A1--A--O};
\foreach \p in {O,O1,A,A1,Q,P} \fill (\p) circle (5pt);
\node[below] at (O){$O$}; \node[above] at (O1){$O_1$}; \node[below] at (A){$A$}; \node[above right] at (A1){$A_1$}; \node[left] at (Q){$Q$}; \node[below] at (P){$P$};
\node[below,font=\small] at ($(P)!0.5!(A)+(0,-0.3)$){$4\sqrt3$};
\node[aux,font=\small] at (5.5,1.9){$R$}; \node[aux,font=\small] at (2.6,7.0){$R$};
\end{tikzpicture}}
Asoslar tomonlari: $\frac{a^2\sqrt3}4=9\sqrt3\Rightarrow a=6$, \ $\frac{b^2\sqrt3}4=81\sqrt3\Rightarrow b=18$.\\
Taşqi çizilgan aylanalar radiuslari: $O_1A_1=\frac6{\sqrt3}=2\sqrt3$, \ $OA=\frac{18}{\sqrt3}=6\sqrt3$. Rasm — öq va $AA_1$ yon qirra orqali kesim.

a) $A_1P\perp OA$: \ $PA=6\sqrt3-2\sqrt3=4\sqrt3$, \ $h=A_1P=PA\cdot\tg60^\circ=12$.
\[V=\frac h3\bigl(S_1+S_2+\sqrt{S_1S_2}\bigr)=4\bigl(9\sqrt3+81\sqrt3+27\sqrt3\bigr)=468\sqrt3.\]
b) Şar markazi $Q$ — $OO_1$ öqda, $OQ=z$. $QA=QA_1=R$ (\textcolor{aux}{qizil}):
\[(6\sqrt3)^2+z^2=(2\sqrt3)^2+(12-z)^2\ \Rightarrow\ 108=12+144-24z\ \Rightarrow\ z=2.\]
\[R^2=108+4=112\ \Rightarrow\ R=4\sqrt7.\]
\Javob{a) $468\sqrt3$; \quad b) $4\sqrt7$}

%================= 45
\newpage
\Savol{45}
\fig{\begin{tikzpicture}[scale=0.8]
\coordinate (P) at (0,0); \coordinate (A) at (4,0); \coordinate (C) at (0,2); \coordinate (T) at (1,1.732);
\fill[okgreen!30] (4,0) arc (0:90:4) -- (0,2) arc (90:0:2) -- cycle;
\draw[thick] (4,0) arc (0:90:4); \draw[thick] (2,0) arc (0:90:2);
\draw[thick] (2,0)--(A); \draw[thick] (0,2)--(0,4);
\draw[dashed,gray] (P)--(2,0); \draw[dashed,gray] (P)--(C); \draw[dashed,gray] (A)--(C);
\draw[axs,line width=1.3pt] (A)--(T); \draw[axs,line width=1.3pt] (T) arc (60:90:2);
\draw[ax] (P)--(T);
\pic[draw,angle radius=0.4cm,"\scriptsize$60^\circ$",angle eccentricity=1.8] {angle=A--P--T};
\foreach \p in {P,A,C,T} \fill (\p) circle (1.6pt);
\node[below left] at (P){$S$}; \node[below] at (A){$A$}; \node[left] at (C){$C$}; \node[above right] at (T){$T$};
\node[below] at (1,0){\small$2$}; \node[below] at (3,0){\small$2$};
\end{tikzpicture}}
Yasovçi $l=\dfrac{R-r}{\cos60^\circ}=\dfrac{2-1}{1/2}=2$.

a) $S_{\text{yon}}=\pi(R+r)l=3\cdot3\cdot2=18$.

b) Kesik konusni töliq konusgaça töldiramiz: uçi $S$, yasovçilari $SA=\frac R{\cos60^\circ}=4$, \ $SC=\frac r{\cos60^\circ}=2$. Yoyilmada sektor burçagi $\varphi=\dfrac{2\pi R}{SA}=\pi$; $A$ va $C$ qarama-qarşi yasovçilarda, ya'ni yoyilmada ular orasidagi burçak $\frac\varphi2=\frac\pi2$ (rasmdagi çorak halqa).\\
$AC$ töğri kesma $S$ ga $\frac{4\cdot2}{\sqrt{20}}\approx1{,}79<2$ masofadan ötadi — sirtdan çiqib ketadi. Şuning uçun eng qisqa yöl: $A$ dan kiçik aylanaga urinma $AT$ va $T$ dan $C$ gaça yoy (\textcolor{aux}{qizil}):
\[AT=\sqrt{4^2-2^2}=2\sqrt3,\quad \cos\angle AST=\frac24\Rightarrow\angle AST=60^\circ,\quad \overset{\frown}{TC}=2\cdot\Bigl(\frac\pi2-\frac\pi3\Bigr)=\frac\pi3\approx1.\]
\[L=2\sqrt3+1.\]
\Javob{a) 18; \quad b) $2\sqrt3+1$}

\ShowJavoblar
\end{document}
